\DTMsavedate{mydate}{2022-06-30}
\maketitle{Electronic Drives}{Electronic Engineering}{\DTMusedate{mydate}}

% ----------------------------------------------------- %
% COURSE INFO
\subsection*{Course Information}
\vspace{0ex}%
\noindent\parbox{0.5\textwidth}{%
    \noindent\begin{tabular}{@{}ll}
                 \textsf{Course:}          & 	Electronic Drives \\
                 \textsf{Course Code:}         & 	TEE 2202    \\
                 \textsf{Credit Hours:}    & 	10
    \end{tabular}}
\vspace{2ex}\hrule\vspace{2ex}

% ----------------------------------------------------- %
% INSTRUCTOR INFO
\subsection*{Lecturer Information}
\noindent\parbox{0.5\textwidth}{%
    \noindent\begin{tabular}{@{}ll}
                 \textsf{Name:}         &    Busiso Mtunzi     \\
% \textsf{Phone:}        &    \phone          \\
\textsf{Email:}        &    busiso.mtunzi@nust.ac.zw  \\
\textsf{Qualifications:}        & PhD in Physics (Renewable Energy), PhD in Physics (PV/T Hybrid Systems),\\
& MSc in Renewable Energy, BTech in Applied Physics
    \end{tabular}}
\vspace{2ex}\hrule\vspace{2ex}


% ----------------------------------------------------- %
% COURSE DESCRIPTION
\subsection*{Course Synopsis}
Power Electronics: Characteristics drive requirements and device protection i.e.\ Converters: DC-DC, DC-AC, AC-AC, AC-DC and control techniques.
Power and distortion factor.
Application of AC and DC motors.
Motor Control i.e.\ variable speed drives, regenerative braking, slip energy recovery, four quadrant operations. Selection and sizing of motor drive systems.
Transducers for power electronics application.
\vspace{2ex} \hrule\vspace{2ex}

% ----------------------------------------------------- %
% LEARNING OBJECTIVES
\subsection*{Learning Objectives}
The course should enable you to:
\begin{outline}
    \1 Understand the operation of power electronic devices in electrical energy control.

    \1 Understand the application of power electronic devices in the lower and upper power end applications of power electronic devices:
    \begin{itemize}
        \item Lower Power End\newline Switched Mode Regulated Power supplies.
        \item Upper Power End\newline These include diverse industrial applications in power electronic converters for variable-speed drives, in steel rolling, textile, paper rolling mills, machine tools and other automation drives.
    \end{itemize}
\end{outline}
\vspace{2ex}\hrule\vspace{2ex}

% ----------------------------------------------------- %
% LEARNING OUTCOMES
\subsection*{Learning Outcomes}
At the conclusion of the course, the student is expected to:
\begin{outline}
    \1     Have a basic understanding of modern power semiconductor devices; their strength, and their switching and protection techniques.
    These include power diodes, bipolar and MOSFET power transistors, silicon controlled rectifiers (SCRs) / thyristors , insulated-gate bipolar transistors (IGBT) and gate turn-off thyristors (GTO).

    \1      Understand the operation and develop analysis skills of DC-DC; DC-AC; AC-DC and AC-AC converters.

    \1      Understand and analyze the qualities of waveforms at input and output ends of these converters (RF; FF; DF).

\end{outline}
\vspace{2ex}\hrule\vspace{2ex}


% ----------------------------------------------------- %
% COURSE TOPICS
\subsection*{Topics}
\begin{outline}
    \1 Introduction\newline Overview of power semiconductor devices, characteristics.
    \1 Diode (Uncontrolled) rectifiers and Harmonic and Distortion factors
    \1 Controlled AC-DC rectifiers on Resistive and Inductive loads
    \1 Choppers, DC-DC converters, Control issues
    \1 DC-AC Converters (Inverters)
    \1 Gate drive and Snubber circuits.
    \1 Device losses and thermal design.
\end{outline}
\vspace{2ex} \hrule\vspace{2ex}

% ----------------------------------------------------- %
% Students Assenssment
\subsection*{Assessment of Students}
There will be two written assignments and two in-class written test (Test 1 on Devices's operation and protection, Test 2 on SCR and Choppers).
At the end of the semester, there shall be a final written examination.
\vspace{2ex}\hrule\vspace{2ex}

% ----------------------------------------------------- %
% Grade Evaluation
\subsection*{Course Evaluation and Grading Policies}
Grades are based on:
Continuous assessment (\qty{25}{\percent}),
Final Exam (\qty{75}{\percent})
\vspace{2ex}\hrule\vspace{2ex}
% ----------------------------------------------------- %
% EQUIPMENT
% https://sites.udel.edu/computing-purchases/personal-specs/
% \subsection*{Essential Equipment}

% \vspace{2ex}\hrule\vspace{2ex}

% Policies and Other information
\subsection*{Policies and Other Information}
% Key Text or some Books
\subsection*{Key Text}
\begin{itemize}
    \item Power Electronics Circuits, Devices and Applications, 3rd edition by M.H. Rashid
    \item Principles of Electric Machines and Power Electronics, 1989 by P. C. Sen
\end{itemize}
\vspace{2ex}\hrule
\vspace{2ex}\hrule\vspace{2ex}